\scanpage{16}
On peut d'ailleurs remplacer la donnée c)
par la donnée
\[
S\in\unsure{\Gamma^0|\Gamma^\circ|\Gamma_0}\backslash E\tag*{c')}
\]
\margin{plus sympathique car on \unsure{passe au|tourne au|prend un|prend une} quotient par un groupe cyclique}
\deleted{\unsure{sont}} et les conditions à mettre sont alors (7) \inserted{et (9)}, puis
\[
\deleted{\Gamma^0\ \text{opère \uline{fid.} sur } S
\quad\left\{\begin{aligned}&\text{ou encor, si }s_0\in S\\&\text{et }u\in\omega\ \text{fixés}\\&us_0\ne s_0\end{aligned}\right.}
\tag*{\deleted{(8')}}
\]
\deleted{\unsure{ou, en av|ou, en ayant}}
\[
\deleted{\unsure{\Gamma|S}\quad\operatorname{card}S=\operatorname{card}(\Gamma^0)=n}\tag*{\deleted{$(8'')$}}
\]
% slanted
\margin{alors \deleted{\unsure{?}}, comme par (9), et $\Gamma\subset\mathit{Aff}(E)$, $\Gamma$ donc $\Gamma^0$ opère \uline{fidèlement} \deleted{dans} \inserted{sur} $S$ on en conclut \deleted{que} (comme $\Gamma^0$ est cyclique et transitif sur $S\ne\varnothing$) que $S$ est un torseur sous $\Gamma^{\unsure{\bullet|0}}$\ldots}
\unsure{puis} Il faut assurer que $S$ est
stable \deleted{sous} $\Gamma$ i.e. est une orbite sous
$\Gamma$, \deleted{Si on choisit \unsure{comme ds|comme dans}} ce qu'on va faire \deleted{en}, \deleted{en
choisissant un} en termes d'un choix
d'un $s_0\in S$, en disant
\[
\exists\,\sigma_{s_0}\in\Gamma,\quad \sigma_{s_0}\ne1,\quad\text{tel que }\sigma s_0=s_0\tag*{(10)}
\]
(NB. A posteriori, ce $\sigma$ sera unique, \deleted{\unsure{et|car}}
par raison combinatoire)

Si on veut \deleted{que} regarder les structures
$P$ dans $E$ entièrement épinglées (choix d'un sommet $s_0$. et de un côté \inserted{comb} $\{s_0,s_1\}$
issu de $s_0$, i.e. d'un sommet $s_1$,
ce qui revient à la donnée de
$u\in\omega$ (\inserted{une} orientation) ($s_1\longleftrightarrow u$ donné par $s_1=us_0$)
% archive sheet number: 15

